\documentclass[11pt]{article}
\usepackage[utf8]{inputenc}
\usepackage{amsmath,amssymb}
\usepackage[margin=1in]{geometry}
\usepackage{hyperref}
\emergencystretch=3em

\title{The log-shift operator:
$\int_0^\infty t^{\nu-1}(c-\log t)^p e^{-\beta t+\beta a\sqrt t}\,dt=(c-\partial_\nu)^p S_\nu(a)$}
\author{Claude, with Daniel Murfet}
\date{2026-09-06}

\begin{document}
\maketitle
\sloppy

\begin{abstract}
Unit 14 of the grammar-paper \S4 autoformalisation: the log-shift operator identity
\texttt{lem:log\_shift}, $\int_0^\infty t^{\nu-1}(c-\log t)^p e^{-\beta t+\beta a\sqrt t}\,dt
=\sum_{m=0}^p\binom{p}{m}c^{p-m}(-1)^m\partial_\nu^m S_\nu(a)$, i.e. $(c-\partial_\nu)^p S_\nu$ written
out. A short binomial argument on top of the general-$k$ order-derivative identity (unit~13). Zero
\texttt{sorry}/\texttt{axiom}.
\end{abstract}

\section{Setting}
Unit~13 gave $\partial_\nu^m S_\nu(a)=\int_0^\infty(\log t)^m t^{\nu-1}e^{\cdots}\,dt$. The log-shift
lemma packages a whole \emph{shifted} logarithm $(c-\log t)^p$ under the integral as the polynomial
$(c-\partial_\nu)^p$ in the order-derivative --- the operator notation the paper uses to compactly write
integrals with log factors.

\section{The thing formalised}
{\small
\begin{verbatim}
theorem fluctuation_log_shift (β a c ν : ℝ) (p : ℕ) (hβ : 0 < β) (hν : 0 < ν) :
    (∫ t in Set.Ioi 0, t^(ν-1) * (c - Real.log t)^p * Real.exp (-β*t + β*a*Real.sqrt t))
      = ∑ m ∈ Finset.range (p+1), (p.choose m : ℝ) * c^(p-m) * (-1)^m
          * iteratedDeriv m (fun l => fluctuation β l a) ν
\end{verbatim}
}

\section{Proof strategy}
Expand $(c-\log t)^p=(-\log t+c)^p=\sum_m(-\log t)^m c^{p-m}\binom{p}{m}$ (\texttt{add\_pow}) pointwise,
distribute the $t^{\nu-1}e^{\cdots}$ factor through the finite sum, and turn $\int\sum$ into $\sum\int$
(\texttt{integral\_finsetSum}, each summand integrable via unit~13's
\texttt{log\_pow\_fluctuation\_integrableOn}). Then \texttt{integral\_const\_mul} pulls the binomial
constants out and \texttt{iteratedDeriv\_fluctuation\_order} identifies each
$\int(\log t)^m t^{\nu-1}e^{\cdots}$ with $\partial_\nu^m S_\nu$.

\section{Roadblocks and resolutions}
None of substance. \texttt{setIntegral\_congr\_fun} again hands an applied-lambda pointwise goal, so
\texttt{beta\_reduce} precedes the binomial rewrite; \texttt{neg\_pow} turns $(-\log t)^m$ into
$(-1)^m(\log t)^m$ and \texttt{ring} matches the per-term factors.

\section{What was Mathlib, what was new}
Mathlib supplied \texttt{add\_pow} (binomial), \texttt{integral\_finsetSum}, \texttt{integral\_const\_mul},
\texttt{neg\_pow}. New is the assembly on top of unit~13: the signed-summand integrability (norm
factorised through \texttt{abs\_pow}) and the finite-sum reorganisation identifying the shifted-log
integral with the order-derivative polynomial.

\section{Lessons}
Once $\partial_\nu^k S=\int(\log)^k\cdots$ is available, any polynomial-in-$\log$ insertion is a
binomial expansion plus \texttt{integral\_finsetSum} --- the only care is proving each signed summand
integrable (bound its norm by the $|\log|^m$-weighted integrand) so the sum/integral swap is licensed.

\section{Follow-ups}
\texttt{lem:log\_generating} ($\sum_p\frac{z^p}{p!}(c-\partial_\nu)^p S_\nu=e^{zc}S_{\nu-z}$) is Taylor's
theorem in the order and needs real-analyticity of $S_\nu$ in $\nu$; \texttt{prop:LogODE} differentiates
the Weber ODE $k$ times in $\lambda$ and needs $\partial_\lambda$/$\partial_a$ commutation. The
ladder-operator commutator $[b,b^\dagger]=1$ still wants a Weyl-algebra framework.
\end{document}
